\documentclass[11pt]{article}
\usepackage[utf8]{inputenc}
\usepackage[english]{babel}
\usepackage{amsmath}
\usepackage{graphicx}
\usepackage{float}
\usepackage{lipsum}
\usepackage{multicol}
\usepackage{xcolor}
\usepackage{tabularx}
\usepackage{booktabs}
\usepackage{hyperref}
\newcolumntype{Y}{>{\centering\arraybackslash}X}
\usepackage[left=2.00cm, right=2.00cm, top=2.00cm, bottom=2.00cm]{geometry}

\title{AN2DL Reports Template}

\begin{document}

\begin{figure}[H]
    \raggedright
    \includegraphics[scale=0.4]{polimi.png} \hfill
    \includegraphics[scale=0.3]{airlab.jpeg}
\end{figure}

\vspace{5mm}

\begin{center}
    % Select between First and Second
    {\Large \textbf{AN2DL - First Homework Report}}\\
    \vspace{2mm}
    % Change with your Team Name
    {\Large \textbf{LosPollosHermanos}}\\
    \vspace{2mm}
    % Team Members Information
    {\large Mohammadhossein Allahakbari,}
    {\large Michele Miotti,}
    {\large Francesco Pesce}\\
    \vspace{2mm}
    % Codabench Nicknames
    {Codabench Nickname1,}
    {michelem,}
    {francescopesce}\\
    \vspace{2mm}
    % Matriculation Numbers
    {Matricola1,}
    {249499,}
    {247974}\\
    \vspace{5mm}
    \today
\end{center}
\vspace{5mm}

\begin{multicols}{2}
    % Note: The following sections represent a suggested
    % structure. We don't need to follow it strictly.

    % -----------------------------------------------------------------------
    % INTRODUCTION
    % -----------------------------------------------------------------------
    \section{Introduction}
    % In this section, you should present your project's context and
    % objectives. You might want to:
    % \begin{itemize}
    %     \item Define the problem (\textit{you may use italics to highlight
    %               definitions})
    %     \item State your goals (\textbf{emphasise key points with bold})
    %     \item Outline your approach
    % \end{itemize}

    % \noindent For instance, you might write: ``This project focuses on
    % \textit{image classification} using \textbf{deep learning} techniques."

    This project presents the results of the \textit{image classification} task
    proposed in the first homework of the Artificial Neural Networks and Deep
    Leaning course at Politecnico di Milano. The goal of the project is to
    classify blood cell images into 8 different classes.

    % -----------------------------------------------------------------------
    % PROBLEM ANALYSIS
    % -----------------------------------------------------------------------
    \section{Problem Analysis}
    % Here you can discuss your initial analysis of the problem. Consider
    % including:
    % \begin{enumerate}
    %     % 8 classes, 96x96 rgb images, labels, etc
    %     \item Dataset characteristics
    %     \item Main challenges % The test set is horrible
    %     That the test set was not horrible? Is this what they mean?
    %     \item Initial assumptions 
    % \end{enumerate}

    % \noindent If you need to reference papers, use the citation command:
    % Recent work~\cite{lecun2015deep} suggests..."

    The dataset consists of 8 classes of blood cell images, each of size 96x96
    pixels, in RGB format. After manually analyzing the dataset, we found that
    a part of it contained duplicate and misleading images that were doctored
    to inject unrelated information into the network.

    Furthermore, the test set provided by the course instructors did not follow
    the same distribution as the training set. This made it difficult to
    evaluate the model's performance accurately by splitting the training set
    into training and validation sets, so we had to rely on the Codabench
    platform for repeated evaluations.

    Overall, the test set's classes were fairly balanced, with each class
    containing a similar number of samples. This made it easier to train the
    model without worrying about class imbalance.

    % -----------------------------------------------------------------------   
    % METHOD
    % -----------------------------------------------------------------------
    \section{Method}
    % Not sure what they are asking here. The final model?
    % This section should detail your approach. You can use equations to
    % explain your methodology. For example, a simple model representation:
    % \begin{equation}
    %     \label{eq:model}
    %     f(x) = \text{softmax}(Wx + b)
    % \end{equation}

    % \noindent Or a more complex loss function:
    % \begin{equation}
    %     \label{eq:loss}
    %     \mathcal{L} = -\frac{1}{N}\sum_{i=1}^{N} y_i\log(\hat{y}_i)
    % \end{equation}

    % \noindent Reference these equations in your text, like:``As shown in
    % equation~\ref{eq:model}..."

    % -----------------------------------------------------------------------
    % EXPERIMENTS
    % -----------------------------------------------------------------------
    \section{Experiments}
    % Step 1: bad CNN just to test things out
    % Step 2: normal augmentation
    % Step 3: overlaying
    % Step 4: keras cv (gigachad picture)
    % Step 5: testing residual networks (up to 0.8)
    %   and inception blocks (up to 0.72)
    % Step 6: custom block (in detail)
    % Step 7: add more filters
    % Step 8.....: not done yet
    % If you need accuracies, number of parameters etc ask me

    % Failed experiments:
    % transfer learning (bad accuracy) (VGG, MobileNet, ResNet), both with and
    % without fine tuning
    % Lion, stochastic gradient descent with momentum

    % Ideas that worked momentarily:
    % voting mechanism with multiple neural networks
    % averaging weights below a certain validation loss
    A number of experiments were conducted to get a higher accuracy on the test
    set. Some of the experiments were successful, while others were not. The
    following is a list of meaningful experiments that were conducted:

    \begin{itemize}
        \item \textbf{Simple CNN:} A simple Convolutional Neural Network (CNN)
              was implemented to test the dataset and the platform. This step
              performed well on our local training and validation sets, but not
              on the Codabench platform. After this step, we suspected that the
              test set was sampled from a different distribution.
        \item \textbf{Data Augmentation:} Normal data augmentation techniques
              were applied to the dataset to increase the number of training
              samples.
        \item \textbf{Overlaying:} Suspecting that the test set was also
              doctored, we overlayed the training set with similar images in
              order to resemble the test set, but this did not improve the
              accuracy in any meaningful way.
        \item \textbf{Keras CV:} The Keras CV library was our first successful
              experiment, achieving an accuracy of 0.71.
        \item \textbf{Residual Networks:} We experimented with Residual
              Networks and Inception Blocks, achieving accuracies of 0.8 and
              0.72, respectively.
        \item \textbf{Custom Block:}
        \item \textbf{More Filters:}
    \end{itemize}

    Some experiments that failed include transfer learning with VGG, MobileNet,
    and ResNet, both with and without fine-tuning. We also tried using a Lion
    optimizer and stochastic gradient descent with momentum. However, these
    experiments did not yield satisfactory results.

    Other ideas were successful at first, such as a voting mechanism with
    multiple neural networks and averaging weights below a certain validation
    loss. These ideas were not pursued further due to their unorthodox nature.
    As the project was a mean to learn from the course, we decided to aproach
    the problem by following the methods taught in class.

    % For your experiments, you might want to present your results in tables.
    % Here's an example of a wide table comparing different models:

    % \begin{table*}[t]
    %     \centering
    %     \setlength{\tabcolsep}{3pt}
    %     \caption{An example of wide table. Best results are highlighted in
    %         \textbf{bold}.}
    %     \begin{tabularx}{\textwidth}{lYYYc}
    %         \toprule
    %         Model            & Accuracy                  & Precision
    %                          & Recall                    & ROC AUC
    %         \\
    %         \midrule
    %         VGG18            & 72.20 $\pm$ 3.06          & 94.95 $\pm$ 0.52
    %                          &
    %         86.95 $\pm$ 0.55 & 80.16 $\pm$ 0.81
    %         \\
    %         Custom Model     & 27.71 $\pm$ 3.19          & 75.70 $\pm$ 1.07
    %                          & 55.75 $\pm$ 2.16          & 36.60 $\pm$ 1.26
    %         \\
    %         ResNet18         & \textbf{89.24 $\pm$ 2.38} & \textbf{95.54
    %         $\pm$ 0.49}      & \textbf{93.43 $\pm$ 1.30} & \textbf{91.68 $\pm$
    %             0.71}
    %         \\
    %         \bottomrule
    %     \end{tabularx}
    %     \label{tab:Performance}
    % \end{table*}

    % \noindent For more specific measurements, you might use a narrower
    % table:

    % \begin{table}[H]
    %     \centering
    %     \setlength{\tabcolsep}{3pt}
    %     \caption{An example of table. Best results may be highlighted in
    %         \textbf{bold}.}
    %     \begin{tabularx}{\linewidth}{lY}
    %         \toprule
    %         Time [$\mu$s] & Distance [mm] \\
    %         \midrule
    %         22$\pm$4      & 8$\pm$1       \\
    %         17$\pm$3      & 7$\pm$1       \\
    %         15$\pm$3      & 6$\pm$1       \\
    %         13$\pm$2      & 5$\pm$1       \\
    %         10$\pm$2      & 4$\pm$1       \\
    %         8$\pm$2       & 3$\pm$1       \\
    %         5$\pm$1       & 2$\pm$1       \\
    %         37$\pm$1      & 1$\pm$1       \\
    %         \bottomrule
    %     \end{tabularx}
    %     \label{tb:Measurements}
    % \end{table}

    % \noindent You can also include figures to visualise your results:
    % \begin{figure}[H]
    %     \centering
    %     \includegraphics[width=0.75\linewidth]{random.jpeg}
    %     \caption{Example figure showing [describe what the figure shows]}
    %     \label{fig:results}
    % \end{figure}

    % \noindent Reference figures using like:``As shown in
    % Figure~\ref{fig:results}..."

    % -----------------------------------------------------------------------
    % RESULTS
    % -----------------------------------------------------------------------
    \section{Results}
    % Keras CV is good
    % Combining residual and inception is good
    % Present your main findings here. You might want to:
    % \begin{itemize}
    %     \item Compare your results with baselines
    %     \item Highlight key achievements using \textbf{bold text}
    %     \item Explain any unexpected outcomes
    % \end{itemize}

    Using the open-source Keras CV library was our most successful experiment,
    as its augmentation pipelines helped us achieve a high accuracy. Trying to
    capture both fine-grained and coarse-grained features, we added inception
    blocks, which improved the accuracy further. Combining residual networks
    and inception blocks, we achieved an accuracy of 0.8.

    % -----------------------------------------------------------------------
    % DISCUSSION
    % -----------------------------------------------------------------------
    \section{Discussion}
    % Wait until we have the final results
    % In this section, analyse your results critically. Consider:
    % \begin{itemize}
    %     \item Strengths and weaknesses
    %     \item Limitations and assumptions
    % \end{itemize}

    % -----------------------------------------------------------------------
    % CONCLUSIONS
    % -----------------------------------------------------------------------
    \section{Conclusions}
    % Summarise your work and discuss potential future directions. This is
    % where you can:
    % \begin{itemize}
    %     \item Restate main contributions
    %     \item Suggest improvements
    %     \item Propose future work
    % \end{itemize}

    Overall, the project's challenge-like nature has helped us follow a
    structured approach to solving the problem. Not expecting a pristine test
    set, we were able to experiment with most of the techniques we learned in
    class, instad of immediately finding the best solution.

    \bibliography{references}
    \bibliographystyle{abbrv}

\end{multicols}
\end{document}